\section{Bilingualism of Canadian Official Languages}\label{bilingualism-of-canadian-official-languages}

\stanceinfo{\textbf{CELC 2024} [2024-01-05]}{PM 2024 [2024-09-22]}

\officialstance{The Canadian Federation of Engineering Students supports Canada's two official languages; French and English as per the Canadian Charter of Rights and Freedoms. As a non-profit organization, the Canadian Federation of Engineering Students holds bilingualism at the highest priority. This is to reduce the systemic gap between the Francophone and Anglophone engineering student membership so that opportunities are not limited because of language barriers.}

\stanceheading{The Students' Position}
\begin{itemize}
 \item There exists a current language barrier between English and French speaking members.
 \item There is a lack of standardized procedures for fully bilingual CFES conferences and activities.
 \item There is a lack of bilingual universities across Canada, as well as a lack of resources to learn a fluent level of either official language within the engineering degree.
 \item The CFES believes all members should feel respected, confident, and heard when speaking or directing any CFES business in either language.
 \item Increasing bilingualism leads to a more inclusive space, with the opportunity to communicate more effectively amongst the membership.
\end{itemize}

\stanceheading{Recommendations to Partners, Stakeholders, and Other Entities}
\begin{itemize}
 \item The CFES calls the Canadian Engineering Accreditation Board (CEAB) to review the language elective policy to provide the option to learn the second official language of Canada, if both are not known.
 \item The CFES calls the CEAB to work alongside Engineering Deans Canada (EDC), to find a place for bilingualism in all schools, regardless of the majority language spoken in an institution.
 \item The CFES calls EDC to support bilingual education and establish it as one of their strategic priorities going forwards in order to close the language gap and provide more opportunities to engineering students on a national level.
 \item The CFES will research partnerships for adult education services that provide French and English education opportunities (tutoring, courses, etc.) to CFES members at a reduced rate. The goal of this is to encourage optional extracurricular learning to bridge the language barrier between Francophone and Anglophone students.
 \item The CFES calls our partners to provide bilingual presentations to ensure that both anglophones and francophones can better understand and follow what is being presented.
 \item The CFES calls the activity managers and CFES officers to ensure all CFES activities remain safe spaces for Francophone minorities through the use of accurate translations and live interpretation.
 \item The CFES calls the members to seek out opportunities to promote bilingualism amongst students and create initiatives that support bilingualism in their institutions.
\end{itemize}

\stanceheading{What the CFES is doing}
\begin{itemize}
 \item Utilizing live interpretation during General Assemblies and other official policy meetings within the CFES.
 \item Investigating the language elective policy within the Canadian Engineering Accreditation Board to see which universities restrict how many courses of the other official language can be taken.
 \item Utilizing current artificial intelligence tools to assist with translation online and in person.
 \item The CFES has a bilingualism commissioner, supervised under the VP Communications, as well as an adjoining Bilingualism Working Group to assist with translation of documents, presentations and policy.
 \item During the election of the CFES national executives, each question must be asked in both official languages by alternating which language is asked first. All candidates will be asked about their level of comprehension of both French and English as well as their plan to communicate with delegates whose first language does not match their own.
\end{itemize}

\stanceheading{What the CFES plans to do}
\begin{itemize}
 \item Create a standardized procedure for translation and interpretation and planning specific procedures for CFES activities to ensure all activities and conferences are fully bilingual and a priority from the beginning.
 \item Provide support to CFES officers who are not fluent in both official languages with resources to help them in the language they are not comfortable in.
 \item Ensure safe spaces are created where both languages can be spoken freely with the support of interpretation when necessary.
 \item Prioritize bilingual presenters in events and activities at CFES conferences, with alternatives for unilingual presenters to be understood by the conference attendees.
 \item Put into practice a Bilingualism Practice Form that provides feedback from members.
 \item Explore partnership with translation services and bilingual focused organizations in the specialization of engineering.
 \item Put into place a system of formal and informal mentorship within the CFES to encourage practicing languages in a low stress environment.
 \item Create a CFES Bilingualism Toolkit for members, that includes but is not limited to the following:
 \begin{itemize}
  \item Bilingualism training
  \item Templates for dual presentation
  \item Best practice for intercommunication
  \item A list of words and expressions
  \item Feedback mechanisms
 \end{itemize}
\end{itemize}

\newpage
